\begin{figure}[htbp]
\centering
\begin{tikzpicture}[x=5.6mm,y=5.6mm]
 \begin{scope}
  \bcuadricula{5}
  \foreach \r in {1,...,5} {\foreach \c in {1,...,5} {
    \pgfmathtruncatemacro{\ataque}{(\r==1)||(\c==2)||(abs(\r-1)==abs(\c-2))}
    \ifnum\ataque=1\relax\bmarca{5}{\r}{\c}\fi
  }}
  \breina{5}{1}{2}
  \foreach \c in {4,5}{\draw[btinta,line width=1.1pt] (\c-.5,3.5) circle[radius=.31];}
  \node[anchor=north,font=\small] at (2.5,-.32) {(a) Padre $B$};
 \end{scope}
 \draw[-{Stealth[length=2mm]},btinta] (5.55,2.5)--(7.35,2.5);
 \node[font=\footnotesize,above] at (6.45,2.65) {$(2,4)$};
 \begin{scope}[shift={(8.5,0)}]
  \bcuadricula{5}
  \foreach \r in {1,...,5}{\foreach \c in {1,...,5}{
    \pgfmathtruncatemacro{\ataque}{(\r==1)||(\c==2)||(abs(\r-1)==abs(\c-2))||(\r==2)||(\c==4)||%
(abs(\r-2)==abs(\c-4))}
    \ifnum\ataque=1\relax\bmarca{5}{\r}{\c}\fi
  }}
  \breina{5}{1}{2}\breina{5}{2}{4}
  \node[anchor=north,font=\small] at (2.5,-.32) {(b) Hijo $B+(2,4)$};
 \end{scope}
\end{tikzpicture}
\caption{La partición exacta del conteo en un tablero pequeño. En la segunda fila del tablero
izquierdo hay dos opciones legales, marcadas con aros. La flecha sigue una de ellas; la otra,
$(2,5)$, da el segundo sumando del conteo. Los puntos señalan casillas ocupadas o atacadas y las
coronas identifican las reinas colocadas. Las casillas disponibles conservan sus coordenadas.}
\label{fig:residual}
\end{figure}
